% GENERATED FILE. Edit canonical lessons, then run npm run generate:books.
% canonical-source-hash: fnv1a64:459f37542f945e6c
% canonical-lessons: GE-C336-einigung, GE-C336-nachgeben, GE-C336-abwaegen, GE-C336-einerseits, GE-C336-andererseits

\chapter{Eine Einigung finden}
\label{ch:ge-eine-einigung-finden}
\hlchaptermodality{336}

\begin{chapteropening}
\textbf{By the end of this chapter:} \emph{I can weigh up both sides of a question, give way and reach an agreement.}

Complete the last lesson of chapter 336: I can weigh up both sides of a question, give way and reach an agreement.
\end{chapteropening}

\section[die Einigung]{die Einigung — the agreement (that is reached)}

\label{lesson:GE-C336-einigung}



\begin{quote}
Before the new one: say the German for to suggest, then the German for to stress.
\end{quote}

\subsection*{You'll want to know: die Einigung}
\textbf{die Einigung} — \textquotedblleft{}the agreement (that is reached)\textquotedblright{}.

\begin{quote}
Nach drei Stunden kamen sie endlich zu einer Einigung.
\end{quote}

\begin{quote}
\emph{After three hours they finally reached an agreement.}
\end{quote}

Plural: \emph{die Einigungen}. It belongs with \emph{sich einigen} — to come to an agreement.

\begin{cousinweb}
\emph{Einigung} is the noun of \emph{sich einigen}, from \emph{einig} \textquotedblleft{}at one, united\textquotedblright{}, built on \emph{ein} \textquotedblleft{}one\textquotedblright{}, inherited from Proto-Germanic \emph{ainaz} (EtymDB): two sides become one.
\end{cousinweb}

\begin{grammarlens}[title={What you've built}]
One more word for telling, explaining and describing at B1.
\end{grammarlens}

\subsection*{Your turn}
\begin{itemize}
\raggedright
  \item \emph{Say it:} \emph{die Einigung}
  \item \emph{Say it:} \emph{die Einigung}, once more
  \item \emph{Say it:} a sentence of your own with \emph{die Einigung}
  \item \emph{Recall:} say the German for to suggest, then the German for to stress, then say \emph{die Einigung} again
\end{itemize}

\begin{tcolorbox}[breakable,colback=teal!4,colframe=teal!35!black,title={Before you move on}]
What does die Einigung mean? (\textquotedblleft{}The agreement (that is reached)\textquotedblright{}.) Say it once more.
\end{tcolorbox}
\section[nachgeben]{nachgeben — to give in, to give way}

\label{lesson:GE-C336-nachgeben}



\begin{quote}
Before the new one: say the German for to stress, then the German for the agreement.
\end{quote}

\subsection*{You'll want to know: nachgeben}
\textbf{nachgeben} — \textquotedblleft{}to give in, to give way\textquotedblright{}.

\begin{quote}
Am Ende hat meine Schwester nachgegeben, und wir haben den Film gesehen.
\end{quote}

\begin{quote}
\emph{In the end my sister gave in, and we watched the film.}
\end{quote}

Separable and strong: \emph{sie gibt nach}, \emph{sie hat nachgegeben}. Also of things: \emph{Der Boden gibt nach} — the ground gives way.

\begin{cousinweb}
\emph{nachgeben} is built from \emph{nach} \textquotedblleft{}after\textquotedblright{} and \emph{geben} \textquotedblleft{}to give\textquotedblright{}, inherited from Old High German \emph{geban} and Proto-Germanic \emph{gebaną}; EtymDB glosses the compound \textquotedblleft{}to give way, to stretch\textquotedblright{}.
\end{cousinweb}

\begin{grammarlens}[title={What you've built}]
One more word for telling, explaining and describing at B1.
\end{grammarlens}

\subsection*{Your turn}
\begin{itemize}
\raggedright
  \item \emph{Say it:} \emph{nachgeben}
  \item \emph{Say it:} \emph{nachgeben}, once more
  \item \emph{Say it:} a sentence of your own with \emph{nachgeben}
  \item \emph{Recall:} say the German for to stress, then the German for the agreement, then say \emph{nachgeben} again
\end{itemize}

\begin{tcolorbox}[breakable,colback=teal!4,colframe=teal!35!black,title={Before you move on}]
What does nachgeben mean? (\textquotedblleft{}To give in, to give way\textquotedblright{}.) Say it once more.
\end{tcolorbox}
\section[abwägen]{abwägen — to weigh up}

\label{lesson:GE-C336-abwaegen}



\begin{quote}
Before the new one: say the German for the agreement, then the German for to give in.
\end{quote}

\subsection*{You'll want to know: abwägen}
\textbf{abwägen} — \textquotedblleft{}to weigh up\textquotedblright{}.

\begin{quote}
Sie wägt die Vorteile und Nachteile genau ab.
\end{quote}

\begin{quote}
\emph{She weighs up the advantages and disadvantages carefully.}
\end{quote}

Separable: \emph{sie wägt ab}. Perfekt: \emph{sie hat abgewogen} (rarely \emph{abgewägt}).

\begin{cousinweb}
\emph{abwägen} is built on \emph{wägen} \textquotedblleft{}to weigh\textquotedblright{}, an older form beside \emph{wiegen}, from Proto-Germanic \emph{weganą} \textquotedblleft{}to move, to carry, to weigh\textquotedblright{}; EtymDB lists English \emph{weigh} and Dutch \emph{wegen} as its cousins.
\end{cousinweb}

\begin{grammarlens}[title={What you've built}]
One more word for telling, explaining and describing at B1.
\end{grammarlens}

\subsection*{Your turn}
\begin{itemize}
\raggedright
  \item \emph{Say it:} \emph{abwägen}
  \item \emph{Say it:} \emph{abwägen}, once more
  \item \emph{Say it:} a sentence of your own with \emph{abwägen}
  \item \emph{Recall:} say the German for the agreement, then the German for to give in, then say \emph{abwägen} again
\end{itemize}

\begin{tcolorbox}[breakable,colback=teal!4,colframe=teal!35!black,title={Before you move on}]
What does abwägen mean? (\textquotedblleft{}To weigh up\textquotedblright{}.) Say it once more.
\end{tcolorbox}
\section[einerseits]{einerseits — on the one hand}

\label{lesson:GE-C336-einerseits}



\begin{quote}
Before the new one: say the German for to give in, then the German for to weigh up.
\end{quote}

\subsection*{You'll want to know: einerseits}
\textbf{einerseits} — \textquotedblleft{}on the one hand\textquotedblright{}.

\begin{quote}
Einerseits ist die Wohnung schön, auf der anderen Seite ist sie sehr teuer.
\end{quote}

\begin{quote}
\emph{On the one hand the flat is nice; on the other, it is very expensive.}
\end{quote}

An adverb: the verb comes right after it, \emph{einerseits ist …}. It usually has a partner later in the sentence, such as \emph{auf der anderen Seite} or \emph{aber}.

\begin{cousinweb}
\emph{einerseits} is built from \emph{ein} \textquotedblleft{}one\textquotedblright{} and \emph{Seite} \textquotedblleft{}side\textquotedblright{}, inherited from Old High German \emph{sīta} and Proto-Germanic \emph{sīdōn} (EtymDB), with the adverb ending \emph{-s}: on one side.
\end{cousinweb}

\begin{grammarlens}[title={What you've built}]
One more word for telling, explaining and describing at B1.
\end{grammarlens}

\subsection*{Your turn}
\begin{itemize}
\raggedright
  \item \emph{Say it:} \emph{einerseits}
  \item \emph{Say it:} \emph{einerseits}, once more
  \item \emph{Say it:} a sentence of your own with \emph{einerseits}
  \item \emph{Recall:} say the German for to give in, then the German for to weigh up, then say \emph{einerseits} again
\end{itemize}

\begin{tcolorbox}[breakable,colback=teal!4,colframe=teal!35!black,title={Before you move on}]
What does einerseits mean? (\textquotedblleft{}On the one hand\textquotedblright{}.) Say it once more.
\end{tcolorbox}
\section[andererseits]{andererseits — on the other hand}

\label{lesson:GE-C336-andererseits}



\begin{quote}
Before the new one: say the German for to weigh up, then the German for on the one hand.
\end{quote}

\subsection*{You'll want to know: andererseits}
\textbf{andererseits} — \textquotedblleft{}on the other hand\textquotedblright{}.

\begin{quote}
Ich möchte gern ein Auto. Andererseits fahre ich gern Fahrrad.
\end{quote}

\begin{quote}
\emph{I would like a car. On the other hand, I like cycling.}
\end{quote}

The partner of \emph{einerseits}, but it can also stand alone. The verb comes right after it: \emph{andererseits fahre ich …}.

\begin{cousinweb}
\emph{andererseits} joins \emph{ander} \textquotedblleft{}other\textquotedblright{} to \emph{Seite} \textquotedblleft{}side\textquotedblright{}, from Old High German \emph{sīta} and Proto-Germanic \emph{sīdōn} (EtymDB), with the adverb ending \emph{-s}.
\end{cousinweb}

\begin{grammarlens}[title={What you've built}]
One more word for telling, explaining and describing at B1.
\end{grammarlens}

\subsection*{Your turn}
\begin{itemize}
\raggedright
  \item \emph{Say it:} \emph{andererseits}
  \item \emph{Say it:} \emph{andererseits}, once more
  \item \emph{Say it:} a sentence of your own with \emph{andererseits}
  \item \emph{Recall:} say the German for to weigh up, then the German for on the one hand, then say \emph{andererseits} again
\end{itemize}

\begin{tcolorbox}[breakable,colback=teal!4,colframe=teal!35!black,title={Before you move on}]
What does andererseits mean? (\textquotedblleft{}On the other hand\textquotedblright{}.) Say it once more.
\end{tcolorbox}
